\section*{الفصل \ref{ch:b2:frontier-orbitals} --- المدارات الحدودية والتفاعلية}

\begin{solution}{exo:b2:frontier-orbitals:1}
الإيثين: \omterm{def:b2:frontier-orbitals:frontier}{HOMO} عند $\alpha + \beta$، \omterm{def:b2:frontier-orbitals:frontier}{وLUMO} عند $\alpha - \beta$. أنيون الأليل (أربعة إلكترونات): \omterm{def:b2:frontier-orbitals:frontier}{HOMO} عند $\alpha$ (المستوى
غير الرابط)، \omterm{def:b2:frontier-orbitals:frontier}{وLUMO} عند $\alpha - 1.414\beta$. البوتاديين: \omterm{def:b2:frontier-orbitals:frontier}{HOMO} عند $\alpha + 0.618\beta$، \omterm{def:b2:frontier-orbitals:frontier}{وLUMO} عند $\alpha -
0.618\beta$.
\end{solution}

\begin{solution}{exo:b2:frontier-orbitals:2}
قاسية: \ce{OH-}، \ce{H+}، \ce{F-}. لينة: \ce{I-}، \ce{RS-}، وكربون اليودوميثان.
\end{solution}

\begin{solution}{exo:b2:frontier-orbitals:3}
يتفاعل البوتاديين (التشكّل s-cis متاح له) والبنتاديين الحلقي (المثبت في s-cis)؛
أما البنتا-1،4-ديين فغير مترافق فلا يتفاعل. \omterm{def:b2:frontier-orbitals:diels-alder}{والديين} (2\textit{Z}،4\textit{Z}) لا يبلغ
الهيئة s-cis إلا بدفع مجموعتي الميثيل الداخليتين إحداهما نحو الأخرى: فهو
يتفاعل بصعوبة شديدة.
\end{solution}

\begin{solution}{exo:b2:frontier-orbitals:4}
سيكلوهكس-3-إين-1-كربونتريل: حلقة سداسية، فيها رابطة مزدوجة بين الكربونين C2
وC3 السابقين في \omterm{def:b2:frontier-orbitals:diels-alder}{الديين}، والمجموعة \ce{C#N} على أحد كربوني الحلقة الناشئين من
\omterm{def:b2:frontier-orbitals:diels-alder}{المحب للديين}.
\end{solution}

\begin{solution}{exo:b2:frontier-orbitals:5}
$\Delta\varepsilon = \qty{8}{eV}$: الاستقرار الاضطرابي $2\beta^2/\Delta\varepsilon = \qty{0.25}{eV}$. والمضبوط: المستوى الأدنى
$-6 - \sqrt{16 + 1} = \qty{-10.123}{eV}$، والربح $2 \times 0.123 = \qty{0.246}{eV}$.
\end{solution}

\begin{solution}{exo:b2:frontier-orbitals:6}
اليودوميثان \omterm{def:b2:frontier-orbitals:hard-soft}{محب للإلكترونات لين}: \omterm{def:b2:frontier-orbitals:control}{تحكم مداري}، فيهاجم بالذرة ذات معامل
\omterm{def:b2:frontier-orbitals:frontier}{HOMO} الأكبر، أي الكربون. والكاتيون الكربوني الحر \omterm{def:b2:frontier-orbitals:hard-soft}{محب للإلكترونات قاسٍ}: \omterm{def:b2:frontier-orbitals:control}{تحكم بالشحنة}،
فيهاجم جزئيًا بالنيتروجين، الذي يحمل شحنة سالبة أكبر.
\end{solution}

\begin{solution}{exo:b2:frontier-orbitals:7}
يرتبط C4 في \omterm{def:b2:frontier-orbitals:diels-alder}{الديين} (أكبر معامل في \omterm{def:b2:frontier-orbitals:frontier}{HOMO}) بالكربون $\beta$ في البروبينال (أكبر
معامل في \omterm{def:b2:frontier-orbitals:frontier}{LUMO}): فينتج 2-ميثوكسي سيكلوهكس-3-إين-1-كربالدهيد، والبديلان على
كربونين متجاورين (ناتج «أورثو»).
\end{solution}

\begin{solution}{exo:b2:frontier-orbitals:8}
هيكل بيسيكلو[2.2.1]هبتين (نوربورنين) تقع فيه مجموعة الإستر على أحد كربوني
\omterm{def:b2:frontier-orbitals:diels-alder}{المحب للديين} السابق، متجهة نحو الرابطة \ce{C=C} الجديدة، تحت الحلقة السداسية
(endo)، في الجهة المقابلة للجسر \ce{CH2}.
\end{solution}

\begin{solution}{exo:b2:frontier-orbitals:9}
تخفض مجموعتا الكربونيل المدار \omterm{def:b2:frontier-orbitals:frontier}{LUMO} \omterm{def:b2:frontier-orbitals:diels-alder}{للمحب للديين}: ففجوته مع \omterm{def:b2:frontier-orbitals:frontier}{HOMO}
البنتاديين الحلقي أصغر بكثير من فجوة الإيثين، والاستقرار $\beta^2/\Delta\varepsilon$
للحالة الانتقالية أكبر، والحاجز أدنى.
\end{solution}

\begin{solution}{exo:b2:frontier-orbitals:10}
التفاعل نوعي فراغيًا: يعطي المالييات ناتج الإضافة الذي تكون فيه مجموعتا الإستر كلتاهما cis
(وكلتاهما endo في الناتج الرئيسي؛ وهو لاكيرالي، مركب meso)؛ ويعطي الفومارات ناتج الإضافة
trans، بإستر endo وآخر exo، متكوّنًا على هيئة زوج من المتماكبات المرآوية (راسيمي).
\end{solution}

\begin{solution}{exo:b2:frontier-orbitals:11}
تصير رابطتا $\pi$ رابطتي $\sigma$ أقوى: $\Delta_r H^\circ < 0$. ويصير جزيئان جزيئًا واحدًا:
$\Delta_r S^\circ < 0$. ويزداد $\Delta_r G^\circ = \Delta_r H^\circ - T\Delta_r S^\circ$ مع $T$ ويصير
موجبًا فوق $T_i = \Delta_r H^\circ/\Delta_r S^\circ$: ففي درجة الحرارة المرتفعة ينشطر ناتج الإضافة عائدًا
(ديلز--ألدر عكسي).
\end{solution}

\begin{solution}{exo:b2:frontier-orbitals:12}
مداران ممتلئان للزوجين الحرين متجاوران يكوّنان تآثرًا رباعي الإلكترونات،
وهو تنافري. فيدور الجزيء حول الرابطة \ce{N-N} حتى يصير الزوجان الحران متعامدين
تقريبًا (تشكّل منحرف)، حيث يكون تداخلهما، ومعه التنافر،
صغيرًا.
\end{solution}

\begin{solution}{pb:b2:frontier-orbitals:1}
\textbf{1.} حلقة خماسية فيها رابطتان مزدوجتان مترافقتان ومجموعة \ce{CH2}: فالحلقة
تثبّت الوحدة الديينية في التشكّل s-cis.
\textbf{2.} \omterm{def:b2:frontier-orbitals:diels-alder}{المحب للديين}، عبر إحدى رابطتيه المزدوجتين.
\textbf{3.} هيكل نوربورنين ملتحم بحلقة بنتين حلقي؛ وتصل رابطتا $\sigma$ الجديدتان
طرفي \omterm{def:b2:frontier-orbitals:diels-alder}{الديين} بكربوني الرابطة المزدوجة \omterm{def:b2:frontier-orbitals:diels-alder}{للمحب للديين}.
\textbf{4.} \omterm{def:b2:frontier-orbitals:endo}{ناتج الإضافة endo} (\omterm{def:b2:frontier-orbitals:endo}{قاعدة endo}): تقع الحلقة الباقية \omterm{def:b2:frontier-orbitals:diels-alder}{للمحب للديين} تحت
\omterm{def:b2:frontier-orbitals:diels-alder}{الديين} في الحالة الانتقالية.
\textbf{5.} نعم: فهو متضافر، وتتكوّن الرابطتان على الوجه نفسه من كل شريك، 
وتُحفظ هندسة \omterm{def:b2:frontier-orbitals:diels-alder}{المحب للديين}.
\textbf{6.} البنتاديين الحلقي \omterm{def:b2:frontier-orbitals:diels-alder}{ديين} جيد (مثبت في s-cis) \omterm{def:b2:frontier-orbitals:diels-alder}{ومحب للديين} في آن واحد، 
\omterm{def:b2:frontier-orbitals:concerted}{والتفاعل متضافر}، بلا مركب وسيط مشحون يحتاج إلى تثبيت.
\textbf{7.} سالبة: تختفي رابطتا $\pi$ وتحل محلهما رابطتا $\sigma$.
\textbf{8.} سالبة: جزيئان يعطيان جزيئًا واحدًا.
\textbf{9.} $\Delta_r G^\circ = \Delta_r H^\circ - T\Delta_r S^\circ$ مع كون الأنتالبي والإنتروبيا كليهما سالبًا: سالبة عند
$T$ المنخفضة، ومعدومة عند $T_i = \Delta_r H^\circ/\Delta_r S^\circ$، وموجبة فوقها.
\textbf{10.} في درجة حرارة الغرفة $T < T_i$: فالتثنّي مفضّل، وإن كان بطيئًا. وقرب
درجة غليان الثنائي يفضّل التوازن الأحادي، ويُبلغ
بسرعة.
\textbf{11.} يغادر الأحادي الخليط الساخن بخارًا: وإزالة ناتج تُبقي
التوازن منزاحًا نحو الأحادي (لوشاتلييه).
\textbf{12.} في درجة حرارة الغرفة يتثنّى من جديد؛ والبرودة تبطّئ ذلك التفاعل.
\textbf{13.} $0.62 - (-0.62) = 1.24|\beta|$.
\textbf{14.} $0.62 - (-0.35) = 0.97|\beta|$.
\textbf{15.} البنتاديين الحلقي مع أنهيدريد الماليك: فالفجوة الأصغر تعطي استقرارًا أكبر
للحالة الانتقالية، وتفاعلًا أسرع.
\textbf{16.} إنهما مترافقتان مع الرابطة \ce{C=C} وساحبتان للإلكترونات: وكما في
البروبينال، تجذبان المستوى $\pi^*$ إلى الأسفل.
\textbf{17.} مداره \omterm{def:b2:frontier-orbitals:frontier}{LUMO}، الذي له فصوص على كربوني الكربونيل مقابل الكربونات المركزية في
مدار \omterm{def:b2:frontier-orbitals:frontier}{HOMO} \omterm{def:b2:frontier-orbitals:diels-alder}{للديين}.
\textbf{18.} هيكل نوربورنين؛ وحلقة الأنهيدريد تحت الرابطة المزدوجة الجديدة، في الجهة المقابلة
للجسر \ce{CH2}.
\textbf{19.} كربونا رأسي الجسر والكربونان الحاملان للأنهيدريد: أربعة
مراكز فراغية.
\textbf{20.} لا: يمر مستوي مرآة عبر الجسر \ce{CH2} وبين نصفي
الجزيء؛ فهو مركب meso.
\textbf{21.} كانا في وضع cis على \omterm{def:b2:frontier-orbitals:diels-alder}{المحب للديين}، والإضافة المتضافرة تبقيهما cis.
\textbf{22.} الهيكل نفسه مع حلقة الأنهيدريد في جهة الجسر \ce{CH2}.
\textbf{23.} لا: فتحويل endo إلى exo يعني كسر الروابط \ce{C-C} وإعادة تكوينها، عبر
التفاعل العكسي، الذي يتطلب التسخين.
\textbf{24.} يفتح الماء الأنهيدريد: فينتج الحمض الثنائي الكربوكسيل cis، ومجموعتا \ce{COOH} كلتاهما endo.
\textbf{25.} \omterm{def:b2:frontier-orbitals:endo}{لناتج الإضافة endo} $\boldsymbol{4}$ مراكز فراغية.
\end{solution}
